% Part: first-order-logic
% Chapter: syntax
% Section: intro-syntax

\documentclass[../../../include/open-logic-section]{subfiles}

\begin{document}

\olfileid{fol}{syn}{itx}

\olsection{పరిచయం}

మొదటిస్థాయి విధేయ తర్కపు సిద్ధాంతాన్నీ అధిసిద్ధాంతాన్నీ అభివృద్ధి
చేయాలంటే, ముందుగా దాని వ్యక్తీకరణల వాక్యనిర్మాణాన్నీ అర్థవిచారాన్నీ
నిర్వచించాలి. మొదటిస్థాయి విధేయ తర్కపు వ్యక్తీకరణలు పదాలు,
\tetoken{సూత్రాలు}{!!{formula}s}. పదాలు \tetoken{చరాలు}{!!{variable}s}, \tetoken{స్థిరసంకేతాలు}{!!{constant}s},
\tetoken{ప్రమేయ సంకేతాలు}{!!{function}s} నుంచి ఏర్పడతాయి. తిరిగి, \tetoken{సూత్రాలను}{!!^{formula}s}
\tetoken{విధేయ సంకేతాలతో}{!!{predicate}s} పాటు పదాల నుంచి (ఇవి అతి చిన్న, ``పరమాణు''
\tetoken{సూత్రాలను}{!!{formula}s} ఏర్పరుస్తాయి), ఆ తరువాత పరమాణు \tetoken{సూత్రాలు}{!!{formula}s} నుంచి
తార్కిక సంయోజకాలు, పరిమాణీకరణికాలను ఉపయోగించి మరింత సంక్లిష్టమైన
వాటిని ఏర్పరుస్తాం. నిర్మాణ నియమాలను పేర్కొనడానికి ఎన్నో పద్ధతులు
ఉన్నాయి; సాధ్యమైన వాటిలో ఒక్కదాన్నే ఇక్కడ ఇస్తాం. ఇతర వ్యవస్థలు వేరే
సంకేతాలను ఎంచుకోవచ్చు, వేరే సంయోజకాల సమితులను ప్రాథమికమైనవిగా
తీసుకోవచ్చు, కుండలీకరణ చిహ్నాలను వేరుగా వాడవచ్చు (లేదా పోలిష్
సంకేతనమని పిలిచే పద్ధతిలోలాగా అసలు వాడకపోవచ్చు). అయితే అన్ని
పద్ధతులకూ ఉమ్మడిగా ఉన్నది ఏమిటంటే, నిర్మాణ నియమాలు పదాలు,
\tetoken{సూత్రాలు}{!!{formula}s} సమితిని \emph{ఆగమనాత్మకంగా} నిర్వచిస్తాయి. సరిగా చేస్తే,
నిర్మాణ నియమాల ప్రకారం ప్రతి వ్యక్తీకరణ సారాంశంగా ఒకే విధంగా
ఉత్పన్నమవుతుంది. వ్యక్తీకరణలు \emph{ఏకార్థంగా పఠించదగినవి} అయ్యేలా
చేసే ఈ ఆగమనాత్మక నిర్వచనం వల్ల, అదే పద్ధతిని---ఆగమనాత్మక
నిర్వచనాన్ని---ఉపయోగించి వాటికి అర్థాలను ఇవ్వగలం.

\end{document}
